\section{Evaluation on NBS Projects}

\label{sec:evaluation}
% ══════════════════════════════════════════════════════════════════════════════

\subsection{Project Portfolio}

We evaluate ZooMRV on 47 nature-based solution projects:

\begin{table}[H]
\centering
\caption{Project portfolio by ecosystem type.}
\label{tab:portfolio}
\begin{tabular}{lcccc}
\toprule
\textbf{Ecosystem} & \textbf{Projects} & \textbf{Area (Mha)} &
  \textbf{Traditional credits} & \textbf{ZooMRV credits} \\
\midrule
Tropical forest & 21 & 1.42 & 18.7M & 12.4M \\
Mangrove & 8 & 0.23 & 2.1M & 1.8M \\
Peatland & 7 & 0.38 & 4.3M & 2.9M \\
Grassland/savanna & 6 & 0.18 & 1.4M & 1.1M \\
Temperate forest & 5 & 0.12 & 0.9M & 0.8M \\
\midrule
\textbf{Total} & \textbf{47} & \textbf{2.33} & \textbf{27.4M} &
  \textbf{19.0M} \\
\bottomrule
\end{tabular}
\end{table}

\subsection{Over-Crediting Analysis}

ZooMRV identifies systematic over-crediting in the traditional portfolio:

\begin{itemize}
  \item \textbf{Inflated baselines}: Dynamic baselines reduce creditable
    deforestation avoidance by 31\% on average, with the largest
    corrections in tropical forest REDD+ projects (37\%).
  \item \textbf{Undetected reversals}: 23\% more reversal events
    detected by continuous monitoring, reducing net credits by 8\%.
  \item \textbf{Biomass overestimation}: CarbonNet corrects systematic
    upward bias in allometric estimates, reducing per-hectare credits by
    12\% in dense tropical forests.
\end{itemize}

Overall, ZooMRV issues 30.7\% fewer credits than traditional
certification for the same project portfolio. This reduction represents
phantom credits that do not correspond to real climate benefit.

\subsection{Uncertainty Quantification}

CarbonNet provides pixel-level uncertainty estimates. At the project
level, we report:

\begin{table}[H]
\centering
\caption{Credit uncertainty by ecosystem type (95\% CI as \% of mean).}
\label{tab:uncertainty}
\begin{tabular}{lcc}
\toprule
\textbf{Ecosystem} & \textbf{ZooMRV uncertainty} &
  \textbf{Traditional uncertainty} \\
\midrule
Tropical forest & $\pm$14\% & $\pm$28\% \\
Mangrove & $\pm$18\% & $\pm$35\% \\
Peatland & $\pm$22\% & $\pm$42\% \\
Grassland & $\pm$19\% & $\pm$38\% \\
\bottomrule
\end{tabular}
\end{table}

ZooMRV reduces credit uncertainty by approximately 50\% across all
ecosystem types, primarily through the fusion of multiple remote sensing
data sources and continuous temporal sampling.


% ══════════════════════════════════════════════════════════════════════════════
